\section{The \texttt{Phonon\_BvK\_PG} McStas Component}
A sample for phonon scattering based on cross section expressions from the Boothroyd textbook
based upon the algorithm from component Phonon\_simple (with expressions from Squires, Ch. 3)
Using the Born-von Karman force constnts and the normal mode formalism

\subsection*{Identification}
\begin{itemize}
  \item \textbf{Author:} Kim Lefmann
  \item \textbf{Origin:} NBI, KU
  \item \textbf{Date:} 06.12.2025
\end{itemize}

\subsection*{Description}
Single-cylinder or slap shaped shape. Absorption included. No multiple scattering. No incoherent scattering emitted, but incoherent is present in attenuation No attenuation from coherent scattering. No Bragg scattering. Specialized for Pyrolytic Graphite (PG)

Algorithm: 0. Always perform the scattering if possible (otherwise ABSORB) 1. Choose direction within a focusing solid angle 2. Select a phonon mode at random (alternatively mode number fixed by user) 3. Calculate the zeros of (E\_i-E\_f-hbar omega(kappa)) as a function of k\_f 4. Choose one value of k\_f (there always at least one possible; see e.g. Squires) 5. Perform the correct weight transformation For details: see article A. F. Davidsen et al, manuscript for J Neutron Res 2026

\subsection*{Input parameters}
Parameters in \textbf{boldface} are required; the others are optional.

\begin{longtable}{p{0.22\textwidth}p{0.12\textwidth}p{0.46\textwidth}p{0.14\textwidth}}
\toprule
\textbf{Name} & \textbf{Unit} & \textbf{Description} & \textbf{Default} \\
\midrule
\endhead
hh & rlu & h-coordinate (in units of a*) of the q-point used when calculating the dispersion & 0 \\
kk & rlu & k-coordinate (in units of a*) of the q-point used when calculating the dispersion & 0 \\
ll & rlu & l-coordinate (in units of c*) of the q-point used when calculating the dispersion & 0 \\
radius & m & Outer radius of sample in (x,z) plane & 0 \\
xwidth & m & Horiz. dimension of sample if slap & 0 \\
yheight & m & Height of sample in y direction & 0 \\
zdepth & m & Depth of sample if slap & 0 \\
sigma\_abs & barns & Absorption cross section at 2200 m/s per atom & 0 \\
sigma\_inc & barns & Incoherent scattering cross section per atom & 0 \\
DW & 1 & Debye-Waller factor; scalar value (no q-dependence) & 1 \\
\textbf{T} & K & Temperature &  \\
focus\_r & m & Radius of sphere containing target. & 0 \\
focus\_xw & m & horiz. dimension of a rectangular area & 0 \\
focus\_yh & m & vert.  dimension of a rectangular area & 0 \\
focus\_aw & deg & horiz. angular dimension of a rectangular area & 0 \\
focus\_ah & deg & vert.  angular dimension of a rectangular area & 0 \\
target\_x & m & position of target to focus at . Transverse coordinate & 0 \\
target\_y & m & position of target to focus at. Vertical coordinate & 0 \\
target\_z & m & position of target to focus at. Straight ahead. & 0 \\
target\_index & 1 & relative index of component to focus at, e.g. next is +1 & 0 \\
\textbf{mode\_input} & 1 & order of the phonon mode to be scattered (unphysical). In the range 0 through 11 . If mode\_input == 12, then the mode is choosen at random &  \\
\textbf{e\_steps\_low} & 1 & Amount of possible intersections beneath the elastic line &  \\
\textbf{e\_steps\_high} & 1 & Amount of possible intersections above the elastic line &  \\
verbose\_input & 1 & toggles verbose mode (verbose\textgreater{}1) or eigenvector mode (verbose==1) & 0 \\
dispersion & 1 & 0: do nothing special; 1: calculate and output the dispersion; 2: calculate and output the dispersion while varying v\_f & 0 \\
\bottomrule
\end{longtable}

\subsection*{Links}
\begin{itemize}
  \item Component source code found in file \texttt{Phonon\_BvK\_PG.comp}.
  \item The test/example instrument \htmladdnormallink{Test\_Phonon.instr}{../examples/Test\_Phonon.instr}.
\end{itemize}
\IfFileExists{contrib/Phonon_BvK_PG_static.tex}{\input{contrib/Phonon_BvK_PG_static.tex}}{}